
\subsubsection{RQ1: Pre-Alignment Margin Predicts Post-Alignment Argmax Flip (H-E1)}

\textbf{Gate result: PASS.}

Table 1 presents the mixed-effects logistic regression results and AUROC values for both model pairs.

\textbf{Table 1. H-E1 existence test results.}

\begin{table*}[htbp]
\centering
\resizebox{\dimexpr 2\columnwidth+\columnsep\relax}{!}{%
\begin{tabular}{llllllllllll}
\toprule
Pair & Alignment & Base model & Aligned model & $\beta _1$ & p-value & AUROC (MMLU) & AUROC (TruthfulQA) & AUROC (ARC) & $\eta^2$ & Flip rate & Gate \\
\midrule
pair2 & DPO & tulu-2-7b & tulu-2-dpo-7b & $-4.3295$ & 4.$13\times 10^{-227}$ & 0.8668 & 0.8034 & 0.9086 & 0.2892 & 12.5\% & PASS \\
pair4 & SFT & pythia-6.9b & oasst-pythia-6.9b & $-0.0617$ & 1.$95\times 10^{-3}$ & 0.6087 & 0.6542 & 0.5756 & 0.0284 & 75.5\% & Below threshold \\
\bottomrule
\end{tabular}
}
\end{table*}

For pair 2 (tulu-2-7b DPO), $\beta _1 = -4.3295$ (p $\approx$ 4.$13\times 10^{-227})$. The negative sign confirms the predicted direction: items with higher pre-alignment confidence margin are less likely to have their argmax changed after DPO alignment. The partial $\eta^2$ = 0.2892 indicates that margin accounts for approximately 29\% of the variance in flip outcomes for pair 2, substantially exceeding the gate threshold of 0.06. Cross-benchmark AUROC values are 0.8668 (MMLU), 0.8034 (TruthfulQA), and 0.9086 (ARC-Challenge), all exceeding the 0.75 threshold without retraining between benchmarks.

For pair 4 (pythia-6.9b SFT), $\beta _1 = -0.0617$ (p = 1.$95\times 10^{-3})$, directionally consistent with pair 2 but substantially weaker. AUROC on MMLU is 0.6087, below the 0.75 threshold. The flip rate for pair 4 is 75.5\%, substantially higher than pair 2's 12.5\%, which may reflect differences in the alignment procedures, model architectures, or training data between the two pairs. Partial $\eta^2$ = 0.0284 for pair 4, below the 0.06 medium-effect threshold.

The cross-benchmark generalization for pair 2 (training on MMLU, evaluation on TruthfulQA and ARC-Challenge without any domain adaptation) is the stronger test of generalizability: the AUROC values of 0.8034 and 0.9086 on held-out benchmark distributions indicate the margin-flip relationship is not limited to the MMLU domain.

\textbf{Empirical flip rate by quintile (pair 2, MMLU).} Items in Q1 (lowest margin) flip at approximately 25\%; items in Q5 (highest margin) flip at approximately 1.5\%. Intermediate quintiles are approximately Q2 $\approx$ 18\%, Q3 $\approx$ 12\%, Q4 $\approx$ 6\%, tracing a monotone decreasing gradient across the confidence spectrum.

Pairs 1 and 3 were excluded due to infrastructure failures (HuggingFace model unavailability and tokenizer errors, respectively). No PPO-aligned model pairs were available.

\begin{figure}[htbp]
\centering
\includegraphics[width=\columnwidth]{figures/fig1_gate_metrics.png}
\caption{Gate metrics for H-E1 --- AUROC and $\beta _1$ per pair}
\label{fig:fig1_gate_metrics}
\end{figure}

\begin{figure}[htbp]
\centering
\includegraphics[width=\columnwidth]{figures/fig3_roc_curves.png}
\caption{ROC curves for pair 2 across MMLU, TruthfulQA, ARC-Challenge}
\label{fig:fig3_roc_curves}
\end{figure}

\begin{figure}[htbp]
\centering
\includegraphics[width=\columnwidth]{figures/fig2_quintile_flip_pair2.png}
\caption{Quintile flip rates for pair 2 (MMLU)}
\label{fig:fig2_quintile_flip_pair2}
\end{figure}


\subsubsection{RQ2: Alignment-Induced Logit Perturbations Are Non-Isotropic (H-M1)}

\textbf{Gate result: PASS.}

Table 2 presents the eigenvalue anisotropy ratios for both evaluated pairs across three benchmarks, alongside the isotropic Gaussian control.

\textbf{Table 2. H-M1 anisotropy results.}

\begin{table*}[htbp]
\centering
\resizebox{\dimexpr 2\columnwidth+\columnsep\relax}{!}{%
\begin{tabular}{llllll}
\toprule
Pair & Alignment & Dataset & Anisotropy ratio ($\lambda _{1}/mean(\lambda _{2},\lambda _{3},\lambda _{4}))$ & p-value (one-tailed) & Significant ($\alpha =0.05)$ \\
\midrule
pair2 & DPO & MMLU (N=14,042) & 2.8996 & 0.0028 & Yes \\
pair2 & DPO & TruthfulQA (N=817) & 3.8281 & 0.0029 & Yes \\
pair2 & DPO & ARC-Challenge (N=1,172) & 2.3360 & 0.0048 & Yes \\
pair4 & SFT & MMLU (N=14,042) & 4.5789 & 0.0047 & Yes \\
pair4 & SFT & TruthfulQA (N=817) & 5.1789 & 0.0053 & Yes \\
pair4 & SFT & ARC-Challenge (N=1,172) & 4.2936 & 0.0072 & Yes \\
Isotropic Gaussian control & --- & N=1,000 & 1.1289 & --- & --- \\
\bottomrule
\end{tabular}
}
\end{table*}

Both evaluated pairs show anisotropy ratios substantially above the isotropic control of 1.13 on all three benchmarks. For pair 2 (DPO), ratios range from 2.34 (ARC) to 3.83 (TruthfulQA) across datasets; the primary MMLU result is 2.90 (p = 0.0028). For pair 4 (SFT), ratios range from 4.29 (ARC) to 5.18 (TruthfulQA); the primary MMLU result is 4.58 (p = 0.0047). The gate criterion ($\geq 2$ of the evaluated model families passing both criteria) is satisfied by both families evaluated.

The MMLU eigenvalue spectra for pair 2 are $\lambda _{1}=7.40, \lambda _{2}=3.26, \lambda _{3}=2.33, \lambda _{4}=2.06$, yielding a ratio of 2.90. For pair 4: $\lambda _{1}=1.22, \lambda _{2}=0.42, \lambda _{3}=0.27, \lambda _{4}=0.10$, yielding a ratio of 4.58. In both cases the dominant eigenvalue is substantially larger than the trailing eigenvalues, indicating that alignment-induced logit perturbations are concentrated along a principal axis in the 4D logit space rather than spread uniformly across directions.

The isotropic sanity check confirms the null model does not spuriously trigger the gate: N(0, $I_{4})$ synthetic data (N=1,000, seed=1) yields ratio = 1.1289.

\begin{figure}[htbp]
\centering
\includegraphics[width=\columnwidth]{figures/fig1_anisotropy_gate_metrics.png}
\caption{Anisotropy ratios per pair vs. isotropic control}
\label{fig:fig1_anisotropy_gate_metrics}
\end{figure}

\begin{figure}[htbp]
\centering
\includegraphics[width=\columnwidth]{figures/fig2_eigenvalue_spectrum.png}
\caption{Eigenvalue spectra $\lambda _1--\lambda _4$ for pair 2 and pair 4}
\label{fig:fig2_eigenvalue_spectrum}
\end{figure}


\subsubsection{RQ3: DPO and SFT Quintile Variance Profiles (H-M2)}

\textbf{Gate result: NULL\_RESULT (LIMITATION\_RECORDED). The directional hypothesis is not supported.}

The H-M2 hypothesis predicted that DPO would produce higher mean logit delta variance in low-confidence regions (Q1) than SFT, after KL-divergence control. The results do not support this directional prediction across any of the three benchmarks.

\textbf{Table 3. H-M2 quintile logit delta variance (MMLU), after KL residualization.}

\begin{table}[htbp]
\centering
\resizebox{\columnwidth}{!}{%
\begin{tabular}{lll}
\toprule
Quintile & DPO variance (pair2, $N\approx 2,800/quintile)$ & SFT variance (pair4, $N\approx 2,800/quintile)$ \\
\midrule
Q1 (lowest margin) & 0.7073 & 0.2229 \\
Q2 & 0.9965 & 0.2250 \\
Q3 & 1.1942 & 0.2544 \\
Q4 & 2.6114 & 0.2939 \\
Q5 (highest margin) & 3.3837 & 0.2815 \\
Q5/Q1 ratio & 4.$79\times$ & 1.$26\times$ \\
\bottomrule
\end{tabular}
}
\end{table}

\textbf{Directional test.} One-tailed Welch's t-test (DPO Q1 mean > SFT Q1 mean): MMLU p = 1.000, Cohen's d = $-0.490$; TruthfulQA p = 1.000, Cohen's d = $-1.536$; ARC-Challenge p = 0.992, Cohen's d = $-0.225$. The direction is reversed across all three benchmarks: SFT mean Q1 delta variance exceeds DPO mean Q1 delta variance.

The H-M2 directional hypothesis is not confirmed. This is recorded as a limitation; the SHOULD\_WORK gate allows the pipeline to continue.

\textbf{Observed quintile trend (post-hoc descriptive finding).} A notable pattern emerges from the quintile profiles: DPO variance increases monotonically from Q1 to Q5 (Q5/Q1 ratio = 4.$79\times$ on MMLU), while SFT variance is nearly flat across quintiles (Q5/Q1 ratio $\approx$ 1.$26\times )$. This profile pattern is consistent across TruthfulQA (DPO: 0.$749\rightarrow 2.807$; SFT: 0.$412\rightarrow 0.751)$ and ARC-Challenge (DPO: 1.$952\rightarrow 3.990$; SFT: 0.$272\rightarrow 0.298)$.

The observed pattern --- DPO applies larger logit perturbations where the base model was already most confident, while SFT distributes perturbations uniformly --- is the opposite of the predicted direction. This finding is descriptive and post-hoc; it is reported as an empirical observation rather than a confirmed mechanistic claim. The finding is also subject to model identity confounds: pair 2 (tulu-2-7b) and pair 4 (pythia-6.9b) differ in architecture, training data, and base model, so the observed difference in quintile profiles cannot be attributed exclusively to the DPO vs. SFT distinction.

\begin{figure}[htbp]
\centering
\includegraphics[width=\columnwidth]{figures/fig1_q1_variance_bar.png}
\caption{Q1 variance comparison: DPO vs SFT per dataset}
\label{fig:fig1_q1_variance_bar}
\end{figure}

\begin{figure}[htbp]
\centering
\includegraphics[width=\columnwidth]{figures/fig2_quintile_trend.png}
\caption{Quintile variance trend lines $Q1\rightarrow Q5$ for DPO and SFT}
\label{fig:fig2_quintile_trend}
\end{figure}

